\section{阶段结论}
\label{cr:conclusion}
本项目已形成AI辅助的DRAM单粒子效应评估流程：限定自然语言请求映射到冻结
参数矩阵，真实Geant4电离沉积经等效源进入TCAD，再由匹配对照、参考支路、
数值检查和来源归档共同支持指标提取。
v1完成了小矩阵扫描及有限外推演示；v2与D1说明读出初态和参考选择会实质改变
结果解释；R50-main在独立物理修订后，完成参考、源隔离、代表点验证及
32个独立训练事件的64个配对响应，并冻结后续64条测试位预测。

四个训练岭模型在工况分组交叉验证中优于均值基准，说明基于沉积量、栅长和
方向的低复杂度代理值得继续验证。当前结论限于这一训练阶段和未校准二维模型；
域内留出、尺寸外推及工艺扩展须按原设计完成真实评分后再加入最终结论。
报告保留失败、未决与未完成状态，使后续更新可以增加证据，而不必改写历史。
课程最终交付还应补齐真实成员贡献、教师指导及作品演示记录。
